\section{Setup}
\label{sec:setup}

\subsection{Task and corpus}
We study step-level grading of mathematical derivations produced by a
retrieval-augmented \citep{lewis2020rag} assistant over a single fixed corpus in
signals and systems. The generator that produced the frozen derivations is
\texttt{qwen-3-235b-a22b-instruct-2507} served via the Cerebras OpenAI-compatible
API (temperature $0$), as recorded in the schema-v2 envelope of the release
eval run. That model ID has since been withdrawn by the provider (HTTP 404);
we therefore treat the derivation set as a fixed artifact and do not regenerate
it under a replacement model.
Each query has a structured expected answer, and the assistant emits an ordered
list of steps, each a symbolic expression together with a short natural-language
description of the operation performed. A grader inspects one step at a time,
given the question and all prior steps as context, and returns one of
\textsc{passed}, \textsc{failed}, or \textsc{abstain}. We map these to
$\{\textsc{true}, \textsc{false}, \textsc{none}\}$ so that every grader is scored
on an identical covered/abstain split (\S\ref{sec:protocol}).

Our benchmark contains 30 queries and is the \emph{source} of the graded
derivation traces, not a separate evaluation set. One query (Q12) failed during
generation and has no label template, so all human-labeled results use the
remaining 29 queries; we report this as a known gap rather than silently
dropping it. Of those 29, 27 come from the primary frozen generator run and two
(Q01, Q07) from an earlier run with the same model ID and temperature $0$; the
released label templates are the ground truth for all grading tables.

\subsection{Symbolic verifier}
The symbolic verifier uses SymPy \citep{meurer2017sympy} to test equivalence
between a step and the derivation context, with additional matchers for composite
expressions and inequalities. By construction it commits only when it can
establish or refute equivalence and abstains otherwise; it therefore never
``guesses.''

\subsection{LLM judges}
The LLM judge is a swappable client selected by environment variables, so a new
judge model flows through the entire pipeline without code changes. Each judge
receives a fixed rubric and step template (reproduced in
Appendix~\ref{app:prompt}) and returns a verdict, a confidence, and a short
rationale. We evaluate two judges chosen from \emph{independent} labs so the
paper can test whether findings are model-specific:
\texttt{gpt-oss-120b} \citep{openai2025gptoss}, an open-weight reasoning model, and
\texttt{claude-sonnet-4-5}, a proprietary model. Using an open-weight judge is a
deliberate reproducibility choice: a reviewer can run the exact grader.

\subsection{Error-injection suite}
To obtain controlled ground truth, we programmatically mutate \emph{9
hand-authored gold derivations} (standard signals/algebra results such as the
Laplace and Fourier transforms of $e^{-at}u(t)$, partial fractions, and the
time-shift property) with 14 mutation operators grouped into five families. The
gold derivations are independent of the generator and verified to pass the
symbolic checker unmutated, so any flag on a mutated step is attributable to the
mutation rather than to a mislabeled baseline. Three families
inject genuine errors: \emph{algebraic} (e.g.\ sign flips), \emph{bounds}
(integration-limit errors), and \emph{identifier} (variable-identity errors).
Two families are negatives that a grader should \emph{not} flag:
\emph{control} (no-op) and \emph{equivalent} (equivalence-preserving rewrites).
This lets us measure both recall on real errors and false positives on
legitimate reformulations. The 9 gold derivations contain 14 mutable step
sites in total; every mutation operator is attempted against every site, for
$14 \times 14 = 196$ attempted (step $\times$ mutation) trials. Of these,
78 are skipped because the mutation does not apply to that step's expression;
the graded suite is the remaining 118 trials (94 error, 24 non-error:
14 control $+$ 10 equivalent).
